% Standalone proof. Shared notation and prose conventions: analyses/README.md.
\documentclass[11pt,a4paper]{article}
\usepackage[T1]{fontenc}
\usepackage[utf8]{inputenc}
\usepackage{amsmath,amsthm}
\usepackage{newtxtext,newtxmath}
\usepackage[a4paper,margin=21mm,headheight=15pt]{geometry}
\usepackage{microtype,booktabs,array,tabularx,longtable}
\usepackage[table]{xcolor}
\usepackage{tikz,fancyhdr,enumitem,seqsplit}
\definecolor{Teal}{HTML}{176B70}
\definecolor{Ink}{HTML}{203443}
\definecolor{Pale}{HTML}{EDF5F5}
\definecolor{Line}{HTML}{C7D8DB}
\definecolor{Muted}{HTML}{536974}
\definecolor{Warm}{HTML}{F8F4EC}
\usepackage[colorlinks=true,linkcolor=Teal,urlcolor=Teal]{hyperref}
\color{Ink}
\newcommand{\blank}{\lambda}
\newcommand{\digits}{\{0,1\}}
\newcommand{\ph}[1]{\mathsf{#1}}
\newcommand{\EXPERT}{\mathcal E}
\newcommand{\CTRL}{\mathcal C}
\newcommand{\PTO}{\operatorname{PTO}}
\newcommand{\triple}[3]{(#1,#2,\mathsf{#3})}
\newcommand{\code}[1]{\texttt{#1}}
\newcommand{\hash}[1]{{\scriptsize\ttfamily\seqsplit{#1}}}
\newtheorem{theorem}{Theorem}
\newtheorem{lemma}{Lemma}
\newtheorem{corollary}{Corollary}
\setlength{\parindent}{0pt}
\setlength{\parskip}{6pt}
\setlist[itemize]{nosep,leftmargin=1.2em}
\pagestyle{fancy}
\fancyhf{}
\fancyhead[L]{\small\sffamily\color{Muted}Binary Multiplication}
\fancyhead[R]{\small\sffamily\color{Muted}Annealing}
\fancyfoot[L]{\scriptsize\sffamily\color{Muted}Exact PTO Equivalence}
\fancyfoot[R]{\small\thepage}
\renewcommand{\headrulewidth}{.35pt}
\renewcommand{\footrulewidth}{0pt}
\newcommand{\sectiontitle}[2]{\refstepcounter{section}\vspace*{2pt}{\LARGE\bfseries\color{Teal}\thesection\quad #2}\par\vspace{7pt}}
\date{}
\pdfinfoomitdate=1
\pdftrailerid{}
\begin{document}
{\Huge\bfseries Exact PTO Equivalence\\[3pt]Binary Multiplication}\par
\vspace{10pt}
{\Large Annealed Controller}\par
{\small\color{Muted}\path{analyses/annealing/DFST_Softmax/mult2/hidden_64_data_100/}}\par
\vspace{7pt}

\vspace{4pt}
\sectiontitle{1 \enspace STATEMENT AND SCOPE}{Statement and Scope}\label{sec:common1}
\textbf{Controllers.}
The exported table \path{minimized_transition_table.pdf} defines controller
$\CTRL$, with initial state $0$ and halt state $51$.
The function \code{mult\_grid(2)} defines expert $\EXPERT$, with initial
state $0$ and halt state $31$. The adjacent expert export matches
this function entry for entry. All tables retain the exported state numbers.

\textbf{Input Domain.}
Let $D=\digits$. Row $0$ initially contains the contiguous word $a\,x\,b$,
starting at column $0$, with $a\in D^m$, $b\in D^n$, and $m,n\ge1$.
All remaining cells contain blanks. The head starts at $(0,0)$.
Write $A=\operatorname{val}_{2}(a)$ and $B=\operatorname{val}_{2}(b)$.
The domain permits arbitrary operand lengths, unequal lengths, and leading zeros.

\textbf{Execution Convention.}
The grid uses $\mathbb Z^2$, with coordinates increasing rightward and
downward. The symbol $\blank$ denotes a blank. A transition
$(q',w,M)$ writes $w$, performs motion $M$, and enters state $q'$.
Motions $L,R,D,U$ move one cell. Motion $S$ writes and stops without moving.

\textbf{Policy-Trajectory Observations.}
For input $a\,x\,b$, define the policy-trajectory observation (PTO) by
\[
\PTO(a\,x\,b)=\bigl((r_t,w_t,M_t)\bigr)_{t=0}^{T-1},
\]
where $r_t$, $w_t$, and $M_t$ denote the read symbol, written symbol,
and motion at controller step $t$. A finite execution contributes the final
$S$ action exactly once, without halt-loop padding.

\begin{theorem}[Exact PTO Equivalence]\label{thm:main}
For every input in the stated domain,
$\PTO_{\CTRL}(a\,x\,b)=\PTO_{\EXPERT}(a\,x\,b)$.
Both executions terminate on the same step. The complete grids and head
positions agree at every step, including the final configuration.
The final grid contains only the $(m+n)$-bit representation of $AB$,
with leading-zero padding, on row $n-1$ in columns $-n,\ldots,m-1$.
The final head stops at $(-n-1,n-1)$.
\end{theorem}

The theorem concerns autonomous executions of the exported discrete controller
on the stated input domain.
\clearpage
\sectiontitle{2 \enspace EXPERT BEHAVIOR}{Expert Behavior}\label{sec:common2}
\subsection{Copying and Partial Products}\label{sec:part1}
Write $b=b_{n-1}\cdots b_0$, so $B=\sum_{i=0}^{n-1}b_i2^i$.
Assign global weight $2^{m-1-z}$ to column $z$.
Column $m-1$ carries the units digit of the final product.

\begin{lemma}[Copy Staircase]\label{lem:copy}
Before product formation, the expert constructs exactly $n$ occupied rows.
For $0\le i<n$, row $i$ contains
\[
a\,x\,b_{n-1}\cdots b_i
\quad\text{starting in column }-i. \tag{1}
\]
The last row contains one right-operand bit. All remaining cells contain blanks.
\end{lemma}
\begin{proof}
Row $0$ contains the input. On row $i$, the expert scans to the first blank,
steps left to $b_i$, then steps left again. A digit at this position implies
at least two remaining right-operand bits. The expert copies every symbol
except $b_i$ one row down and one column left, proceeding from right to left.
Every destination contains a blank before writing; the source word remains unchanged.

The first descent reaches a blank on the new row. Each later descent reaches
the copied right neighbor: a digit, except directly below the last digit of
$a$, where the copied separator $x$ occupies the lower cell. The cell below
the source separator contains a digit. After copying the first digit of $a$,
the expert reads the blank immediately to the left and descends to the first
digit of the new row.

The new row satisfies (1) for $i+1$, with one fewer right-operand bit.
When one right-operand bit remains, the second left step reads $x$ and starts
product formation. Thus $n-1$ copies produce exactly $n$ occupied rows.
\end{proof}

\begin{lemma}[Aligned Partial Products]\label{lem:products}
After product formation, row $i$ contains the $(m+1)$-bit representation of
$A b_i$ in columns $-i-1,\ldots,m-i-1$, with a leading $0$.
In global column weights, row $i$ represents $A b_i2^i$.
All remaining cells contain blanks.
\end{lemma}
\begin{proof}
The expert forms products from row $n-1$ upwards. On row $i$, the expert
erases the rightmost factor bit $b_i$, the remaining factor prefix, and $x$,
then walks left across $a$. For $b_i=0$, the expert writes zeros; for $b_i=1$,
the expert preserves the digits of $a$. Both branches write one additional
$0$ on the left blank. The resulting word has the claimed support and value,
with a leading zero because $A b_i<2^m$.

The expert scans right to the blank at column $m-i$ and moves up.
For $i>0$, the upper cell contains the last digit of $a$ on row $i-1$;
the next cell contains $x$. The expert scans the remaining factor prefix,
selects the last bit, and repeats. Each upper prefix contains at least two
bits, so another digit follows the selected bit during leftward erasure.
For $i=0$, the upper cell contains a blank, ending product formation and
starting summation on row $0$.
\end{proof}

Both passes terminate: each copy shortens the remaining factor prefix, and
product formation visits each of the $n$ rows once. Every horizontal scan
crosses a finite word.

\newpage
\subsection{A Fixed-Width Addition Invariant}\label{sec:part2}
\begin{lemma}[Accumulated Product]\label{lem:merge}
After merging through row $k$, $0\le k<n$, all rows above $k$ contain only blanks.
Row $k$ contains a contiguous $(m+k+1)$-bit accumulator on columns
$-k-1,\ldots,m-1$, with value
\[
S_k=A\sum_{i=0}^{k}b_i2^i.
\]
Rows below $k$ retain the partial products from Lemma~\ref{lem:products}.
At the inspection of the next row, the expert occupies state $20$ with head
$(-k-2,k+1)$.
\end{lemma}
\begin{proof}
For $k=0$, the initial left scan crosses the product on row $0$, reads the
leading $0$, reaches the blank at column $-2$, and descends.

If row $k+1$ exists, the inspected cell contains the leading $0$ at column
$-k-2$. The expert scans right across the lower product to the first blank,
at column $m-k-1$. Moving up reaches an accumulator digit because
\[
-k-1\ \le\ m-k-1\ \le\ m-1.
\]
The expert scans the remaining upper digits, reaches the right blank, and
steps left to column $m-1$. The lower cell contains a blank because the
original lower product ends at $m-k-2$. The first column therefore adds
the upper bit to zero without a carry.

The expert processes the remaining columns from right to left. For upper
bit $u$, lower bit $v$ (with a blank counting as zero), and incoming carry
$c$, the update satisfies
\[
u+v+c=w+2c',\qquad w\in\{0,1\},\quad c'\in\{0,1\}. \tag{2}
\]
The expert erases the upper bit, writes $w$ below, and stores $c'$ in the
control state. The processed lower suffix contains the correct sum suffix,
the unprocessed cells retain the previous values, and the carry records
the contribution to the next column. Equation~(2) preserves this invariant.

The original lower word extends one column farther left than the upper word,
with a leading $0$ at $-k-2$. After consuming the upper word, carry zero
preserves this bit, whereas carry one changes the bit to $1$ and clears.
The expert reaches the blank at $-k-3$ on row $k+1$ and descends.
Row $k$ now contains only blanks, and row $k+1$ represents
$S_k+A b_{k+1}2^{k+1}=S_{k+1}$ on the claimed support.
\end{proof}

\begin{lemma}[Carry Confinement]\label{lem:carry}
During an upper-digit update with incoming carry $1$, the lower cell
contains an original binary digit. If the carry outlives the upper word,
the next lower cell contains the original leading $0$.
\end{lemma}
\begin{proof}
A blank gap lies to the right of the original lower product. The first
column has carry zero, and $u+0+0\le1$, so no carry originates in the gap.
A carry first arises inside the lower digit block. Further leftward steps
remain inside that block until the leading zero. The upper word ends one
column before that zero, where $0+0+1=1$ clears the carry.
\end{proof}

Each merge consumes finitely many upper digits and at most one extra carry
column. After $n-1$ merges, state $20$ reads the blank below the final row
and moves up to state $30$, which emits the final $S$ action and enters halt
state $31$. Lemma~\ref{lem:merge} gives the $(m+n)$-bit product and final head
$(-n-1,n-1)$. All other cells contain blanks.

\clearpage
\sectiontitle{3 \enspace OBSERVATION PROTOCOL}{Observation Protocol}\label{sec:common3}
The observation protocol refines expert states by geometric context. Each phase records the permitted read symbols and successor phases. Section~\ref{sec:common4} gives the complete protocol together with the controller-state relation; the shared tables avoid repeating the same phase definitions.

\begin{lemma}[Protocol Coverage]\label{lem:coverage}
Every expert execution on a well-formed multiplication follows the
$94$ phases in Section~\ref{sec:common4}, from phase $0$ to phase $93$.
\end{lemma}
\begin{proof}
The copy and product invariants (Lemmas~\ref{lem:copy}--\ref{lem:products})
give all copying and product-reading conditions. The scan annotations
remember unchanged cells and distinguish one digit from at least two;
the explanations in Sections~\ref{sec:part3}--\ref{sec:part5} prove these refinements.
Lemmas~\ref{lem:merge}--\ref{lem:carry} and the blank-gap argument in
Section~\ref{sec:part6} supply the addition and boundary conditions. Induction over
each transition, each copied row and each merged row establishes coverage.
The finite-pass arguments in the expert invariants give termination in phase $93$.
\end{proof}

\clearpage
\sectiontitle{4 \enspace SIMULATION CERTIFICATE}{Simulation Certificate}\label{sec:common4}
\subsection{Copying and Remembered Factor Bits}\label{sec:part3}
The next four tables define a finite protocol and simulation relation.
Phase $i$ pairs expert state $g_i$ with a controller state $q\in Q_i$,
with identical grids and heads. An entry $x:(j,w,M)$ means read $x$,
write $w$, move $M$, and advance to \emph{phase} $j$.
The symbol $d\in\{0,1\}$ means either digit, consistently within that entry.
For \emph{each} $q\in Q_i$, the actual successor lies in $Q_j$.

\begin{center}
\fontsize{9.2}{11}\selectfont\renewcommand{\arraystretch}{1.21}
\setlength{\tabcolsep}{4pt}
\begin{tabularx}{\textwidth}{r p{5.5cm} c c X}
\toprule\rowcolor{Pale}$i$ & Phase & $g_i$ & $Q_i$ & Read $:$ (next phase, write, move)\\\midrule
0 & First left-operand digit & 0 & $\{0\}$ & $d:\triple{1}{d}{R}$ \\
1 & Scan left operand & 0 & $\{0,24\}$ & $d:\triple{1}{d}{R}\quad x:\triple{2}{x}{R}$ \\
2 & First right-operand digit & 0 & $\{7\}$ & $0:\triple{3}{0}{R}\quad 1:\triple{4}{1}{R}$ \\
3 & Scan factor; last bit 0 & 0 & $\{48\}$ & $0:\triple{3}{0}{R}\quad 1:\triple{4}{1}{R}\quad \blank:\triple{5}{\blank}{L}$ \\
4 & Scan factor; last bit 1 & 0 & $\{48\}$ & $0:\triple{3}{0}{R}\quad 1:\triple{4}{1}{R}\quad \blank:\triple{6}{\blank}{L}$ \\
5 & Revisit last factor bit 0 & 1 & $\{14\}$ & $0:\triple{7}{0}{L}$ \\
6 & Revisit last factor bit 1 & 1 & $\{14\}$ & $1:\triple{8}{1}{L}$ \\
7 & Copy or use remembered bit 0 & 2 & $\{49\}$ & $0:\triple{9}{0}{D}\quad 1:\triple{11}{1}{D}\quad x:\triple{25}{x}{R}$ \\
8 & Copy or use remembered bit 1 & 2 & $\{19\}$ & $0:\triple{9}{0}{D}\quad 1:\triple{11}{1}{D}\quad x:\triple{26}{x}{R}$ \\
9 & First copied 0: descend & 3 & $\{16\}$ & $\blank:\triple{10}{\blank}{L}$ \\
10 & First copied 0: write & 6 & $\{13\}$ & $\blank:\triple{13}{0}{U}$ \\
11 & First copied 1: descend & 4 & $\{4\}$ & $\blank:\triple{12}{\blank}{L}$ \\
12 & First copied 1: write & 7 & $\{41\}$ & $\blank:\triple{13}{1}{U}$ \\
13 & Continue after first copy & 9 & $\{35\}$ & $0:\triple{14}{0}{D}\quad 1:\triple{16}{1}{D}\quad x:\triple{19}{x}{D}$ \\
14 & Copy 0 above a lower digit & 3 & $\{44,47\}$ & $d:\triple{15}{d}{L}$ \\
15 & Write copied 0 & 6 & $\{9,49\}$ & $\blank:\triple{18}{0}{U}$ \\
16 & Copy 1 above a lower digit & 4 & $\{26\}$ & $d:\triple{17}{d}{L}$ \\
17 & Write copied 1 & 7 & $\{2,3\}$ & $\blank:\triple{18}{1}{U}$ \\
18 & Continue the copy & 9 & $\{34\}$ & $0:\triple{14}{0}{D}\quad 1:\triple{16}{1}{D}\quad x:\triple{19}{x}{D}\quad \blank:\triple{24}{\blank}{D}$ \\
19 & Descend below operator & 5 & $\{45\}$ & $d:\triple{20}{d}{L}$ \\
20 & Write shifted operator & 8 & $\{5\}$ & $\blank:\triple{21}{x}{U}$ \\
21 & Copy last left-operand digit & 9 & $\{34\}$ & $0:\triple{22}{0}{D}\quad 1:\triple{23}{1}{D}$ \\
22 & Copied 0 above lower operator & 3 & $\{44\}$ & $x:\triple{15}{x}{L}$ \\
23 & Copied 1 above lower operator & 4 & $\{26\}$ & $x:\triple{17}{x}{L}$ \\
24 & Begin scan on copied row & 0 & $\{36\}$ & $d:\triple{1}{d}{R}$ \\
\bottomrule
\end{tabularx}
\end{center}

\textbf{Coverage.} Both operands contain at least one digit. The initial scan therefore
meets a digit, $x$, and a right-operand digit before a blank. Phases $3$--$8$
remember the last factor bit: the scan's last digit, the digit revisited
by the first left step, and (if only one factor bit remains) the digit
selected for multiplication all refer to the same unchanged cell.

Lemma~\ref{lem:copy} proves the remaining restrictions: the first descent
sees a blank; subsequent descents see the already copied right-neighbor
digit, except below the last left-operand digit, where the lower cell
contains $x$. Every destination contains a blank. Copying $x$ exposes a
source digit. Restarting on a copied row also reads a digit. These conditions
hold for arbitrary operand lengths and values.

\newpage
\subsection{Forming and Scanning Partial Products}\label{sec:part4}
The product scan starts just to the right of the new padding zero.
``One'' denotes a scan after exactly one product digit; the remaining
scan phases follow at least two product digits. Each scan also records the most recently read digit.

\begin{center}
\fontsize{9.2}{11}\selectfont\renewcommand{\arraystretch}{1.21}
\setlength{\tabcolsep}{4pt}
\begin{tabularx}{\textwidth}{r p{5.5cm} c c X}
\toprule\rowcolor{Pale}$i$ & Phase & $g_i$ & $Q_i$ & Read $:$ (next phase, write, move)\\\midrule
25 & Select remembered factor bit 0 & 10 & $\{10\}$ & $0:\triple{27}{\blank}{L}$ \\
26 & Select remembered factor bit 1 & 10 & $\{3\}$ & $1:\triple{28}{\blank}{L}$ \\
27 & Erase operator for factor 0 & 11 & $\{10\}$ & $x:\triple{29}{\blank}{L}$ \\
28 & Erase operator for factor 1 & 12 & $\{4\}$ & $x:\triple{31}{\blank}{L}$ \\
29 & First digit of product by 0 & 13 & $\{25\}$ & $d:\triple{30}{0}{L}$ \\
30 & Continue product by 0 & 13 & $\{25\}$ & $d:\triple{30}{0}{L}\quad \blank:\triple{33}{0}{R}$ \\
31 & First digit of product by 1 & 14 & $\{6\}$ & $d:\triple{32}{d}{L}$ \\
32 & Continue product by 1 & 14 & $\{6\}$ & $d:\triple{32}{d}{L}\quad \blank:\triple{33}{0}{R}$ \\
33 & First digit after product padding & 15 & $\{15,36\}$ & $0:\triple{34}{0}{R}\quad 1:\triple{35}{1}{R}$ \\
34 & After one product digit 0 & 15 & $\{0\}$ & $0:\triple{36}{0}{R}\quad 1:\triple{37}{1}{R}\quad \blank:\triple{38}{\blank}{U}$ \\
35 & After one product digit 1 & 15 & $\{24\}$ & $0:\triple{36}{0}{R}\quad 1:\triple{37}{1}{R}\quad \blank:\triple{39}{\blank}{U}$ \\
36 & Product scan; last digit 0 & 15 & $\{24\}$ & $0:\triple{36}{0}{R}\quad 1:\triple{37}{1}{R}\quad \blank:\triple{40}{\blank}{U}$ \\
37 & Product scan; last digit 1 & 15 & $\{24\}$ & $0:\triple{36}{0}{R}\quad 1:\triple{37}{1}{R}\quad \blank:\triple{41}{\blank}{U}$ \\
38 & Inspect above one-digit product 0 & 16 & $\{39\}$ & $d:\triple{42}{d}{R}\quad \blank:\triple{52}{\blank}{D}$ \\
39 & Inspect above one-digit product 1 & 16 & $\{17\}$ & $d:\triple{42}{d}{R}\quad \blank:\triple{53}{\blank}{D}$ \\
40 & Inspect above; last digit 0 & 16 & $\{17\}$ & $d:\triple{42}{d}{R}\quad \blank:\triple{54}{\blank}{D}$ \\
41 & Inspect above; last digit 1 & 16 & $\{17\}$ & $d:\triple{42}{d}{R}\quad \blank:\triple{55}{\blank}{D}$ \\
42 & Cross upper operator & 17 & $\{8,43\}$ & $x:\triple{43}{x}{R}$ \\
43 & First bit in upper factor prefix & 17 & $\{16\}$ & $0:\triple{44}{0}{R}\quad 1:\triple{45}{1}{R}$ \\
44 & Scan factor; last bit 0 & 17 & $\{16\}$ & $0:\triple{44}{0}{R}\quad 1:\triple{45}{1}{R}\quad \blank:\triple{46}{\blank}{L}$ \\
45 & Scan factor; last bit 1 & 17 & $\{16\}$ & $0:\triple{44}{0}{R}\quad 1:\triple{45}{1}{R}\quad \blank:\triple{47}{\blank}{L}$ \\
46 & Select remembered factor bit 0 & 10 & $\{13\}$ & $0:\triple{48}{\blank}{L}$ \\
47 & Select remembered factor bit 1 & 10 & $\{13\}$ & $1:\triple{50}{\blank}{L}$ \\
48 & First remaining bit to erase (0) & 11 & $\{27\}$ & $d:\triple{49}{\blank}{L}$ \\
49 & Erase factor prefix (0) & 11 & $\{10\}$ & $d:\triple{49}{\blank}{L}\quad x:\triple{29}{\blank}{L}$ \\
50 & First remaining bit to erase (1) & 12 & $\{22\}$ & $d:\triple{51}{\blank}{L}$ \\
51 & Erase factor prefix (1) & 12 & $\{1\}$ & $d:\triple{51}{\blank}{L}\quad x:\triple{31}{\blank}{L}$ \\
\bottomrule
\end{tabularx}
\end{center}

\textbf{Coverage.} Lemma~\ref{lem:products} gives a nonempty block of $m$
product digits after the padding zero. This justifies the first-digit phase
and the split into exactly one or at least two scanned digits. At the right blank, the most recently read digit equals the units digit.

The bottom row contains one factor bit; erasing that bit exposes $x$.
Every higher row contains at least two factor bits, so another digit
precedes $x$ during erasure. An upward return from a completed row reaches
either a blank or the last left-operand digit followed by $x$.
The upper factor scan records the last bit, which remains unchanged until selection. All these facts follow from the copy staircase.

\newpage
\subsection{Starting the Sum and Aligning a Row}\label{sec:part5}
These phases retain the last product digit until the next visit and
distinguish a one-digit left operand from a longer operand. The remembered
bit and length class serve as proof annotations without adding machine registers.

\begin{center}
\fontsize{9.2}{11}\selectfont\renewcommand{\arraystretch}{1.21}
\setlength{\tabcolsep}{4pt}
\begin{tabularx}{\textwidth}{r p{5.5cm} c c X}
\toprule\rowcolor{Pale}$i$ & Phase & $g_i$ & $Q_i$ & Read $:$ (next phase, write, move)\\\midrule
52 & Begin one-digit sum; last bit 0 & 18 & $\{20\}$ & $\blank:\triple{56}{\blank}{L}$ \\
53 & Begin one-digit sum; last bit 1 & 18 & $\{42\}$ & $\blank:\triple{57}{\blank}{L}$ \\
54 & Begin sum; remember last digit 0 & 18 & $\{42\}$ & $\blank:\triple{58}{\blank}{L}$ \\
55 & Begin sum; remember last digit 1 & 18 & $\{42\}$ & $\blank:\triple{59}{\blank}{L}$ \\
56 & Revisit sole product digit 0 & 19 & $\{28\}$ & $0:\triple{60}{0}{L}$ \\
57 & Revisit sole product digit 1 & 19 & $\{11\}$ & $1:\triple{60}{1}{L}$ \\
58 & Revisit product units digit 0 & 19 & $\{11\}$ & $0:\triple{62}{0}{L}$ \\
59 & Revisit product units digit 1 & 19 & $\{11\}$ & $1:\triple{62}{1}{L}$ \\
60 & Read the sole leading zero & 19 & $\{2\}$ & $0:\triple{61}{0}{L}$ \\
61 & Left boundary after sole digit & 19 & $\{23\}$ & $\blank:\triple{66}{\blank}{D}$ \\
62 & Second digit of initial left scan & 19 & $\{2,26\}$ & $d:\triple{63}{d}{L}$ \\
63 & Third digit of initial left scan & 19 & $\{2,23\}$ & $0:\triple{64}{0}{L}\quad 1:\triple{65}{1}{L}$ \\
64 & Left scan after reading zero & 19 & $\{23\}$ & $0:\triple{64}{0}{L}\quad 1:\triple{65}{1}{L}\quad \blank:\triple{66}{\blank}{D}$ \\
65 & Left scan after reading one & 19 & $\{23\}$ & $0:\triple{64}{0}{L}\quad 1:\triple{65}{1}{L}$ \\
66 & Inspect fresh lower leading cell & 20 & $\{12,36\}$ & $0:\triple{67}{0}{R}\quad \blank:\triple{92}{\blank}{U}$ \\
67 & First digit after lower leading zero & 21 & $\{0\}$ & $d:\triple{68}{d}{R}$ \\
68 & Scan the lower partial product & 21 & $\{24\}$ & $d:\triple{68}{d}{R}\quad \blank:\triple{69}{\blank}{U}$ \\
69 & First upper alignment digit & 28 & $\{17\}$ & $d:\triple{70}{d}{R}$ \\
70 & Next upper alignment cell & 28 & $\{8\}$ & $d:\triple{71}{d}{R}\quad \blank:\triple{72}{\blank}{L}$ \\
71 & Finish aligning the upper row & 28 & $\{4,8\}$ & $d:\triple{71}{d}{R}\quad \blank:\triple{72}{\blank}{L}$ \\
\bottomrule
\end{tabularx}
\end{center}

\textbf{Initial Left-Scan Coverage.}
On the top row, the product's units digit has not changed between the right scan and the first left step of summation. If $m=1$, the full padded word equals $0d$: scanning left reads $d$, then $0$, then a blank (phases
$56$--$57$, $60$--$61$). If $m\ge2$, the padded word has at least three digits,
so two more digits must follow the first one (phases $62$--$63$).
In both cases the padding zero immediately precedes the left blank.

\textbf{Alignment Coverage.}
Lemma~\ref{lem:merge} gives either a blank in the inspected lower cell
(no remaining row) or the next row's leading $0$. After that zero comes
at least one digit. The upward move from the lower right boundary reaches
an upper digit, so alignment begins inside the upper word. The expert scans to the right blank and then returns left to the units digit. The first lower units cell contains a blank. These facts justify every alignment branch here and
the first three phases in the next certificate subsection.

\newpage
\subsection{Column Addition, Carries and Stopping}\label{sec:part6}
Separate phases track carry-zero returns from lower blanks and lower digits. This distinction prevents an
impossible early end of the upper word from entering the proof protocol.

\begin{center}
\fontsize{9.2}{11}\selectfont\renewcommand{\arraystretch}{1.21}
\setlength{\tabcolsep}{4pt}
\begin{tabularx}{\textwidth}{r p{5.5cm} c c X}
\toprule\rowcolor{Pale}$i$ & Phase & $g_i$ & $Q_i$ & Read $:$ (next phase, write, move)\\\midrule
72 & Read upper units digit & 22 & $\{32,41\}$ & $0:\triple{73}{\blank}{D}\quad 1:\triple{74}{\blank}{D}$ \\
73 & Write units digit 0 below & 24 & $\{27,46\}$ & $\blank:\triple{88}{0}{U}$ \\
74 & Write units digit 1 below & 25 & $\{22\}$ & $\blank:\triple{88}{1}{U}$ \\
75 & Return above, carry zero & 23 & $\{50\}$ & $\blank:\triple{76}{\blank}{L}$ \\
76 & Read upper bit, carry zero & 22 & $\{37\}$ & $0:\triple{77}{\blank}{D}\quad 1:\triple{78}{\blank}{D}\quad \blank:\triple{83}{\blank}{D}$ \\
77 & Add value zero below & 24 & $\{30\}$ & $d:\triple{75}{d}{U}\quad \blank:\triple{90}{0}{U}$ \\
78 & Add value one below & 25 & $\{29\}$ & $\blank:\triple{90}{1}{U}\quad 0:\triple{75}{1}{U}\quad 1:\triple{79}{0}{U}$ \\
79 & Return above, carry one & 27 & $\{10\}$ & $\blank:\triple{80}{\blank}{L}$ \\
80 & Read upper bit, carry one & 26 & $\{38\}$ & $0:\triple{81}{\blank}{D}\quad 1:\triple{82}{\blank}{D}\quad \blank:\triple{85}{\blank}{D}$ \\
81 & Add one inside lower digit block & 25 & $\{29\}$ & $0:\triple{75}{1}{U}\quad 1:\triple{79}{0}{U}$ \\
82 & Add two inside lower digit block & 29 & $\{9\}$ & $d:\triple{79}{d}{U}$ \\
83 & Preserve unneeded leading zero & 19 & $\{20\}$ & $0:\triple{84}{0}{L}$ \\
84 & Blank to the left of merged row & 19 & $\{31,34\}$ & $\blank:\triple{66}{\blank}{D}$ \\
85 & Apply final carry to leading zero & 25 & $\{42\}$ & $0:\triple{86}{1}{U}$ \\
86 & Return after final carry & 23 & $\{47\}$ & $\blank:\triple{87}{\blank}{L}$ \\
87 & Step down left of merged row & 22 & $\{33\}$ & $\blank:\triple{84}{\blank}{D}$ \\
88 & Return after units digit & 23 & $\{18,21\}$ & $\blank:\triple{89}{\blank}{L}$ \\
89 & Read the second upper digit & 22 & $\{21\}$ & $0:\triple{77}{\blank}{D}\quad 1:\triple{78}{\blank}{D}$ \\
90 & Return from the lower blank gap & 23 & $\{21\}$ & $\blank:\triple{91}{\blank}{L}$ \\
91 & Upper digit after lower blank gap & 22 & $\{21\}$ & $0:\triple{77}{\blank}{D}\quad 1:\triple{78}{\blank}{D}$ \\
92 & Emit final stop & 30 & $\{35,40\}$ & $\blank:\triple{93}{\blank}{S}$ \\
93 & Halted & 31 & $\{51\}$ & $\text{--}$ \\
\bottomrule
\end{tabularx}
\end{center}

\textbf{The Lower Blank Gap.}
While merging row $k$ into row $k+1$, the upper word begins at $-k-1$;
the original lower word ends at $m-k-2\ge-k-1$. A blank below an upper digit
must therefore lie to the right of that lower endpoint. The next column
left still has index at least $-k-1$, so the next upper reading contains a digit.
This proves the restriction in phase $91$. Independently, the upper word
has length $m+k+1\ge2$, proving the second-digit restriction in phase $89$.

\textbf{Carry and Boundary Coverage.}
Lemma~\ref{lem:carry} excludes lower blanks in the carry-one digit phases
and forces $0$ for the final carry extension. Returning above a written
lower digit reads the upper blank just erased. When the upper word ends,
the next lower cell contains the original leading zero: carry zero
preserves this bit, whereas carry one changes the bit to $1$. Both branches reach the blank to the
left of the merged row. When no next row remains, the upward move reaches
a blank, and the expert emits the final $S$. Phase $93$ has no further PTO actions.

\textbf{Local Transition Identity.}
For every listed input branch $x:(j,w,M)$ and every $q\in Q_i$, direct
comparison with the two exported tables gives
\begin{equation}
\delta_\EXPERT(g_i,x)=(g_j,w,M),\qquad
\delta_\CTRL(q,x)=(q',w,M)\quad\text{for some }q'\in Q_j. \tag{3}
\end{equation}
The certificate contains $109$ phase--state pairs and $201$ local checks,
covering $133$ distinct nonterminal controller entries. A transition reaches
a halt state exactly when $j=93$, with $M=S$. Initially $g_0=0$ and $Q_0=\{0\}$.
Protocol coverage excludes the remaining controller entries on valid inputs.

\clearpage
\sectiontitle{5 \enspace EQUIVALENCE PROOF}{Equivalence Proof}\label{sec:common5}
\label{sec:equivalence}
\begin{proof}[Proof of Theorem~\ref{thm:main}]
Initially both configurations have the same grid and head,
with control states $0$ related by phase $0$. Suppose the relation holds
before a controller step. Lemma~\ref{lem:coverage} supplies an allowed read symbol
$x$. Equal grids and heads give the same read symbol to both controllers.
Equation~(3) gives equal writes and motions and preserves the relation
between successor states. The next grids and heads therefore agree.

Induction establishes equality of every finite execution prefix. The terminal
condition permits a halt exactly on an edge to the terminal phase.
The finite copy, product, and merge passes establish termination.
Both controllers therefore halt together with equal complete PTOs and final
configurations. Lemma~\ref{lem:merge} gives the product, padding, and final head.
\end{proof}

\bigskip
\sectiontitle{6 \enspace VERIFICATION AND ARTIFACTS}{Verification and Artifacts}\label{sec:common6}
\textbf{Verification.}
Run \code{python3 verify\_equivalence.py {-}{-}quick} from this proof directory to check
file fingerprints, CSV/JSON agreement, the exported expert against the
\code{mult\_grid(2)} source, and all $201$ certificate obligations.
Omitting \code{{-}{-}quick} additionally compares independent executions on
all $15{,}876$ operand pairs of lengths $1$--$6$, plus $45$ directed and
$128$ fixed-seed random cases, including $1024$-bit operands.
All $10{,}184{,}741$ PTO steps agree. The checks also compare final grids,
heads, and padded products. These tests exercise $132$ distinct nonterminal entries.
The unbounded result follows from the invariants and induction above.

\textbf{Artifact Identity.}
The title page gives the repository-relative case directory. This proof directory
contains the controller and expert exports, reference code, protocol, checker,
and saved verification report. The verification script stores or reads SHA-256 checksums for these artifacts. A checksum identifies the exact bytes
of a file; the checker compares the stored and computed values to detect changes.
These checks bind the certificate to the supplied exports. The invariants,
protocol coverage, and local transition identity establish equivalence.
\end{document}
